\documentclass[12pt]{article}
\usepackage{graphicx,amssymb,amsmath,amsthm}
\usepackage{hyperref,url}
\usepackage{verbatim}
\usepackage[labelfont=bf]{caption}
\newcommand{\nocopyright}{
No Copyright\thanks{
The authors hereby waive all copyright
and related or neighboring rights to this work,
and dedicate it to the public domain.
This applies worldwide.
}}
\title{Taut fillings}
\author{Peter Doyle \and Matthew Ellison \and Zili Wang\thanks{
Zili Wang is supported by
the National Natural Science Foundation of China (Grant No. 12301432)
}}
\date{Version 1.1 dated 28 May 2025
}
\newtheorem{theorem}{Theorem}
\newtheorem{prop}{Proposition}
\newtheorem{lemma}[theorem]{Lemma}
\newtheorem*{hypo}{Hypothesis}
\newtheorem*{question}{Question}
\newtheorem{corollary}{Corollary}
\newtheorem*{conjecture}{Conjecture}

\def\startproof{{\bf {\medskip}{\noindent}Proof: }}
\def\definition{{\bf {\bigskip}{\noindent}Definition: }}
\def\endproof{$\spadesuit$  \newline}

\newcommand{\proofstart}{{\noindent \bf Proof.\ }}
\newcommand{\note}{{\noindent \bf Note.\ }}
\newcommand{\pfend}{\spadesuit}
\newcommand{\mathproofend}{\quad \qed}
\newcommand{\proofend}{$\quad \qed$}
\newcommand{\goesto}{\rightarrow}
\newcommand{\p}{\;\;\;}
\newcommand{\comma}{\p,}
\newcommand{\pd}{\p.}
\newcommand{\semi}{\p;}
\newcommand{\Rplus}{{\RR_{\geq 0}}}



\newcommand{\beginfigurepos}{\begin{figure}[htbp]}
%
% \fig
% followed by {filename}{label}{caption}
%
\newcommand{\fig}[3]{
\beginfigurepos
\includegraphics[width=370pt]{figures/#1.pdf}
\caption{#3}
\label{#2}
\end{figure}
}
%
% \figsize
% followed by {size}{filename}{label}{caption}
%
\newcommand{\figsize}[4]{
\beginfigurepos
\centerline{
\includegraphics[width=#1]{figures/#2.pdf}
}
\caption{#4}
\label{#3}
\end{figure}
}


\def\ZZ{\mathbf{Z}}
\def\QQ{\mathbf{Q}}
\def\RR{\mathbf{R}}
\def\heartsuit{\mbox{\boldmath{$Z$}}}% 
\def\del{\partial}
\def\dist{\mathrm{dist}}
\def\eps{\epsilon}


\newcommand{\parsec}{\subsection*}
\newcommand{\sect}[1]{\S\ #1}

%\newcommand{\gversus}[2]{\braket{#1\,|\,#2}}
\newcommand{\pair}[1]{\braket{#1}}
\newcommand{\sub}[1]{\braket{#1}}
\newcommand{\nnn}{\mathbf{NSet}}
\newcommand{\zzz}{\mathbf{ZSet}}
\newcommand{\nn}{\mathbf{N}}
\newcommand{\zz}{\mathbf{Z}}

\newcommand{\versus}[2]{\left(#1\,|\,#2\right)}
\newcommand{\gversus}[2]{\left[#1\,|\,#2\right]}
\newcommand{\poly}[1]{\left[#1\right]}

\newcommand{\rising}[1]{\overline{#1}}
\newcommand{\floor}[1]{\left \lfloor #1 \right \rfloor}
\newcommand{\ceiling}[1]{\left \lceil #1 \right \rceil}
\newcommand{\gbinom}{\genfrac{[}{]}{0pt}{}}

\makeatletter
\newcommand{\xRrightarrow}[2][]{\ext@arrow 0359\Rrightarrowfill@{#1}{#2}}
\newcommand{\Rrightarrowfill@}{\arrowfill@\equiv\equiv\Rrightarrow}
\newcommand{\xLleftarrow}[2][]{\ext@arrow 3095\Lleftarrowfill@{#1}{#2}}
\newcommand{\Lleftarrowfill@}{\arrowfill@\Lleftarrow\equiv\equiv}
\newcommand{\xLleftRrightarrow}[2][]{\ext@arrow 3399\LleftRrightarrowfill@{#1}{#2}}
\newcommand{\LleftRrightarrowfill@}{\arrowfill@\Lleftarrow\equiv\Rrightarrow}
\makeatother

\newcommand{\eq}{\Rrightarrow}
\newcommand{\funeq}[1]{\stackrel{#1}\eq}
\newcommand{\funequiv}[1]{\stackrel{#1}\equiv}
\newcommand{\funto}[1]{\stackrel{#1}\to}
\newcommand{\union}{\cup}
\newcommand{\cross}{\times}
\newcommand{\op}{\operatorname}
\newcommand{\kw}{\mathrm}
\newcommand{\escape}{\op{escape}}
\newcommand{\cancel}{\op{cancel}}
\newcommand{\nestuntil}{\operatorname{nestuntil}}
\newcommand{\size}{\operatorname{size}}
\newcommand{\id}{\operatorname{id}}
\newcommand{\create}{\operatorname{create}}
\newcommand{\destroy}{\operatorname{destroy}}
\newcommand{\then}{\triangleleft}
\newcommand{\lp}{\kw{LP}}
\newcommand{\LPvol}{\kw{LPvol}}
\newcommand{\vol}{\kw{vol}}
\newcommand{\flipdist}{\kw{flipdist}}
\newcommand{\tetbar}{\kw{homvol}}
\newcommand{\vsa}{\kw{vsa}}
\newcommand{\cone}{\kw{cone}}
\newcommand{\link}{\kw{link}}
\newcommand{\maxdeg}{\kw{maxdeg}}
\newcommand{\mindeg}{\kw{mindeg}}
\newcommand{\vertices}{\kw{vertices}}
\newcommand{\carrier}{\kw{carrier}}
\newcommand{\nbhd}{\kw{nbhd}}
\newcommand{\adj}{\kw{adj}}
\newcommand{\taut}{\kw{taut}}
\newcommand{\net}{\kw{net}}

\newcommand{\Zvol}{\kw{Zvol}}
\newcommand{\Qvol}{\kw{Qvol}}
\newcommand{\Rvol}{\kw{Rvol}}
\newcommand{\tetvol}{\kw{tetvol}}
\newcommand{\hash}{\#}
\newcommand{\ms}{[}
\newcommand{\msend}{]}

\newcommand{\kwstyle}{\bf}
\newcommand{\IF}{{\ \kwstyle{if}\ }}
\newcommand{\ELSE}{{\ \kwstyle{else}\ }}
\newcommand{\WHILE}{{\kwstyle{while}\ }}
\newcommand{\DO}{{\ \kwstyle{do}\ }}
\newcommand{\RETURN}{{\kwstyle{return}\ }}
\newcommand{\direct}{\oplus}




\begin{document}

\maketitle

\begin{abstract}
Let $\sigma$ be a simplicial triangulation of the 2-sphere,
$X$ the associated integral 2-cycle.
A filling of $X$ is an integral 3-chain $Y$
with $\del Y = X$;
a taut filling is one with minimal $L_1$-norm.
We show that any taut filling 
arises from an extension of $\sigma$ to
a simplicial complex homeomorphic to the 3-ball.
This complex is shellable, and
every clique in its 1-skeleton occurs as a simplex.
The key to the proof is the general
fact that any taut filling of an $n$-cycle
splits under disjoint union, connected sum, and more generally
what we call almost disjoint union,
where summands are supported on sets that overlap in at most $n+1$ vertices.
Despite the generality of this result, we have nothing to say about
optimal fillings of spheres of dimension 3 or higher.
\end{abstract}

\section{Overview}
Let $\Delta$ be the abstract $|V|-1$-simplex with vertices $V$,
viewed as a simplicial complex.
Let $C_n$ and $Z_n$ be its integral $n$-chains and $n$-cycles.
Here and throughout we'll take $n \geq 1$.
A \emph{filling} of an $n$-cycle $X \in Z_n$ is any
$n+1$-chain $Y \in C_{n+1}$ with $\del Y = X$.
Let $\Zvol(X)$ be the minimum $L_1$-norm of a filling of $X$,
and call $Y$ \emph{taut} if $|Y|=\Zvol(\del Y)$.

For $X \in C_n$ let
$\vertices(X)$ be the set of all the vertices of all the $n$-simplices
to which $X$ assigns non-$0$ weight.
For a taut filling $Y$ of $X$ we have $\vertices(Y) = \vertices(X)$,
because projecting $V$ onto $\vertices(X)$ by mapping
$V \setminus \vertices(X)$ to an arbitrary $x \in \vertices(X)$
will push any filling with an internal vertex to a smaller filling.
(Cf.\ Proposition \ref{nointernal} below.)
This is why we write $C_n$ and $Z_n$ without reference to $V$.

Call an $n$-cycle $X=X_1+X_2 \in Z_n$ an \emph{almost disjoint union} if
\[
|\vertices(X_1) \cap \vertices(X_2)| \leq n+1
.
\]
This notion generalizes both disjoint union and connected sum
along an $n$-simplex.
Ellison \cite{ellison:gap}
showed the $\Zvol$ adds under almost disjoint union:
$\Zvol(X) = \Zvol(X_1) + \Zvol(X_2)$.
This means that \emph{some} taut filling $Y$ of $X$ splits into a sum
$Y=Y_1+Y_2$ of taut fillings of $X_1,X_2$.
Here we amplify Ellison's Corollary 2 to show (Theorem \ref{th1})
that when $n \geq 2$, \emph{any} taut filling of $X$ splits.

In place of integral chains we can use chains with
coefficients in $\QQ$ or $\RR$.
Evidently $\Rvol=\Qvol$.
Computing $\Rvol$ is a linear programming problem with rational coefficients,
so $\Qvol = \Rvol$.
Ellison used LP duality to prove that
$\Qvol$ adds under almost
disjoint union.
Now for $n \geq 2$ we get the stronger result
(Corollary \ref{cor1})
that taut $\QQ$-fillings split:
Multiplying a taut
$\QQ$-filling $Y$ of $X$ by a common denominator $q$
yields a taut $\ZZ$-filling $qY$ of $qX$:
Split the $\ZZ$-filling $qY$ and divide by $q$ to get a splitting
of the $\QQ$-filling $Y$.

Now let $\sigma$ be a simplicial triangulation of $S^2$;
let $X(\sigma) \in Z_2$ be either one of the 2-cycles
that arises from $\sigma$ by orienting its $2$-simplices;
and write $\Zvol(\sigma) = \Zvol(X(\sigma))$.
Let $\tetvol(\sigma)$ be the minimum number of $3$-simplices
required to extend $\sigma$ to a simplicial triangulation $\tau$ of $B^3$.
We will show that
\[
\Zvol(\sigma) = \tetvol(\sigma)
.
\]
Certainly $\Zvol \leq \tetvol$, because any extension $\tau$ induces a filling
of $X(\sigma)$.
Theorem \ref{th2} states that 
any taut filling of $X(\sigma)$ arises in this way.

The proof of Theorem \ref{th2} relies on Theorem \ref{th1}.
We use induction, with base case $\sigma$ a tetrahedron.
In any taut filling $Y$ of $X(\sigma)$,
thought of as a multiset of oriented $3$-simplices,
some $t \in Y$ meets $\sigma$ in at least two faces.
Ignoring the base case,
$Y-t$ is a taut filling of
$X(\sigma) - \del t = X(\sigma^\prime)$,
where $\sigma^\prime$ is either a simplicial triangulation of $S^2$,
or two such triangulations $\sigma_1,\sigma_2$ fused along an edge.
In the first case induction yields a triangulation of $B^3$,
and we glue on $t$.
In the second case by Theorem \ref{th1}
$Y-t$ splits into taut fillings of $X(\sigma_1),X(\sigma_2)$;
by induction both give triangulations of $B^3$, which we
glue onto $t$.
So we wind up with some kind of triangulation $\tau$ of $B^3$,
built from the simplices of $Y$.
The wrinkle is that we need to ensure that in $Y$ there are no
identifications beyond those in $\tau$.

We don't know what happens for spheres of dimension 3 or greater.

\section{Background}

Our interest in $\Zvol$ stems from the work
of Sleator, Tarjan, and Thurston
\cite{stt:jams}.
They observed that coning from a vertex of maximal degree
shows that a $v$-vertex
triangulation $\sigma$
of $S^2$ has $\tetvol \leq 2v-10$.
They then
produced examples for which $\tetvol = 2v-10$, provided
$v$ is larger than some unspecified bound,
and conjectured that such examples should exist for any $v \geq 13$.
Their approach was to 
realize $\sigma$ as an ideal hyperbolic polyhedron
with volume $2v-O(\log(v))$ when measured in bushels,
a \emph{bushel} being volume of an equilateral ideal tetrahedron,
which is maximal.
This implies $\tetvol \geq 2v-O(\log(v))$.
From there they worked their way up to showing
that $\tetvol=2v-10$.
They suggested
\cite[p. 697]{stt:jams}
that in fact $\Qvol= 2v-10$,
which would immediately imply $\tetvol=2v-10$.
Mathieu and Thurston \cite{mt}
produced a different class of examples
with $\Qvol = 2v-10$, still under the assumption that $v$
is sufficiently large.
In \cite{dew:filling} we produced examples for any $v \geq 13$ with
$\Qvol=2v-10$, confirming the conjecture of Sleator, Tarjan, and Thurston.
Theorem \ref{th2} here now
tells us that whenever we have $\Qvol=2v-10$, any taut filling
must arise from a triangulation of the ball.

About $\Qvol$ versus $\Zvol$.
In \cite{dew:filling}
we describe triangulations of $S^2$ with $\Qvol < \Zvol$.
Since $\Qvol$ and $\Zvol$
add under connected sum, the gap between them can be arbitrarily
large
(Ellison \cite{ellison:gap}).
Some examples with a gap have $\maxdeg =6$,
but in all such examples that we have seen,
the gap is $<1$, so that
$\lceil \Qvol \rceil = \Zvol$.
Taking connected sums of these examples
produces vertices of degree $\geq 7$.


\section{Taut fillings}

As long as $n>0$,
as we will continue to assume throughout,
we can think of an $n$-chain $X \in C_n$ as a multiset of
non-cancelling oriented simplices:
\[
X = \sum_{t \in X} t
.
\]
In this sum $t$ denotes an oriented simplex
\[
t=[x_0,\ldots,x_n]
= [x_{\pi(0)},\ldots,x_{\pi(n)}]
,\;\;\mbox{$\pi$ an even permutation}
,
\]
and each unoriented simplex contributes a number of terms corresponding
to its multiplicity.

The size of $X$ as a multiset is its $L_1$-norm $|X|$.
Write $U \subset X$
if $U$ is a sub-multiset of $X$.
This happens just when
$|X| = |U| + |X-U|$.

\begin{prop} \label{subtaut}
If $Y$ is taut and $U \subset Y$ then $U$ is taut.
\end{prop}

\proofstart
This is clear, but let's drag it out.
If $U$ is not taut, there is a smaller filling $V$ of $\del U$.
Take $Y^\prime = Y-U+V$.
In terms of multisets, this means that we substitute $V$ for $U$,
and then do any required cancellation of oppositely oriented simplices
of $Y-U$ and $V$.
We have
\[
\del Y^\prime = \del Y - \del U + \del V = \del Y
\]
and
\[
|Y^\prime| \leq |Y-U| + |V| = |Y| - |U| + |V| < |Y|
,
\]
contradiction.
\proofend

\section{Coning}
For $x \in V$, $U \in C_n$ let
\[
\nbhd(x,U) = \sum_{ t \in U: x \in \vertices(t) } t
,
\]
\[
\deg(x,U) = |\nbhd(x,U)|
,
\]
\[
\maxdeg(U) = \max_x \deg(x,U)
.
\]
The \emph{cone from $x$ to $U$}
is the $(n+1)$-chain
\[
\cone(x,U) = 
\sum_{t \in U} \adj(x,t)
,
\]
consisting of all the non-trivial oriented
$(n+1)$-simplices $\adj(x,t)=[x x_0\ldots x_n]$ obtained
by adjoining $x$ to  $t = [x_0\ldots x_n] \in U$.
If $x \in \vertices(t)$ then $t$ doesn't contribute to the sum,
so $|\cone(x,U)| = |U| - \deg(x,U)$.

If $X$ is closed then
$\del\, \cone(x,X) = X$, so

\begin{prop} \label{maxdeg}
\[
\Zvol(X) \leq |X| - \maxdeg(X)
\mathproofend
\]
\end{prop}

If $U \in Z_n$ and $x \notin \vertices(U)$ then
$\deg(x,U)=0$ and
$|\cone(x,U)| = |U|$.
In this case we call $\cone(x,U)$ a \emph{complete cone}.
By Proposition \ref{maxdeg} a non-trivial complete cone is not taut.
(This is a long-winded way of saying that instead of coning from
$x \notin U$, we could have coned from some $x \in U$.)
Thus:
\begin{prop}
If $Y$ is taut it contains no non-trivial complete cone.
\proofend
\end{prop}

Call $x \in \vertices(Y) \setminus \vertices(\del Y)$ an
\emph{internal vertex} of $Y \in C_{n+1}$.
If $x$ is internal to $Y$ then $\nbhd(x,Y)$ is a complete cone, so
as observed above:
\begin{prop} \label{nointernal}
If $Y$ is taut then
it has no internal vertices.
\proofend
\end{prop}

\section{Almost disjoint unions}

Recall that if $X,Y \in Z_n$ we call $X+Y$ an almost disjoint union
if
\[
|\vertices(X) \cap \vertices(Y)| \leq n+1
.\]
The most interesting special case is a connected sum,
where $t=[c_0,c_1,\ldots,c_n]$ occurs once in $X$ and
$-t=[c_1,c_0,\ldots,c_n]$ occurs once in $Y$.
For example, if $n=1$ with $X$ a cycle of length $p$
and $Y$ a cycle of length $q$
the connected sum $X+Y$ is
a cycle of length $p+q-2$.
$\Zvol$ adds, because
$\Zvol(X)=p-2$,
$\Zvol(Y) = q-2$,
and
\[
\Zvol(X+Y) = (p+q-2)-2 = \Zvol(X)+\Zvol(Y)
.
\]
In this case not every filling of $X+Y$ splits.
We want to show that fillings do split when $n \geq 2$.

For $A \subset V$, $p \in A$ define
\[
\pi_{A,p}: V \to A
\]
\[
\pi_{A,p}(x) = x \IF x \in A \ELSE p
\]
and let
\[
K_*(A,p): C_*(V) \to C_*(A)
\]
be the induced chain map.
This is a projection of $C_*(V)$ onto $C_*(A)$.

Let $A,B$ be finite vertex sets sharing the vertices
$C = A \cap B$.
For practice with this overloading of the letter `C',
observe that $C_{|C|}(C)$ is trivial, as is $Z_{|C|+1}(C)$.

Let
\[
C_*(A,B) = C_*(A) \direct C_*(B)
.
\]
For $p,q \in C$ define
\[
g_*(p,q) = K_*(A,p) \direct K_*(B,q):
C_*(A \cup B) \to C_*(A,B)
.
\]

\begin{prop} \label{recover}
Let $(X,Y) \in C_n(A,B)$.
If $|C| \leq n$
we can recover $(X,Y)$ from $X+Y$.
If $(X,Y) \in Z_n(A,B)$ this holds also when $|C|  = n+1$.
\end{prop}

\proofstart
We can assume $|C| \geq 1$.
(Add a brand new point to $C$ if necessary.)
We claim that
for any $p,q \in C$
(not necessarily distinct) we have
\[
g_n(p,q)(X+Y)=(X,Y)
.
\]
Indeed,
\[
K_n(A,p)(X+Y) = X + K_n(A,p)(Y)
.
\]
The second term belongs to $C_n(C)$,
which is trivial when $|C| \leq n$.
If $Y \in Z_n(B)$ the second term belongs to $Z_n(C)$,
which is trivial when $|C| \leq n+1$.
\proofend


\begin{theorem} \label{th1}
If $|C| \leq n+1$,
for all $(X,Y) \in Z_n(A,B)$ we have
\[
\Zvol(X+Y) = \Zvol(X)+\Zvol(Y)
.
\]
And as long as $n \geq 2$,
for any $Z \in \taut(X+Y)$ 
we have $Z = Z_X + Z_Y$
with $Z_X \in \taut(X)$, $Z_Y \in \taut(Y)$.
\end{theorem}

\note
This result resembles Theorem 6.2 of
Pournin and Wang
\cite{pw:flip}
about flip paths.
Like theirs, our proof uses a variation on the normalization technique
of STT
\cite[Lemma 7]{stt:jams}.
It would be nice to fit these results under one roof.
\medskip

\proofstart
We can assume $|C| = n+1$,
as this is the hardest case.
And we might as well go ahead and take $n=2$, $|C|=3$,
as this case illustrates all the issues.

Take any $Z \in \taut(X+Y) \subset C_3(A \union B)$.
Pick distinct points $p,q \in C$,
and let
\[
(Z_X,Z_Y) = g_{n+1}(p,q)(Z)
.
\]
(For now we'll suppress the dependence of $Z_X,Z_Y$ on $p,q$.)

Because $g_*(p,q)$ is a chain map we have
$\del_{n+1} Z_X = X$, $\del_{n+1} Z_Y = Y$:
\[
(\del_{n+1} Z_X, \del_{n+1} Z_Y)
=
\del_{n+1} (g_{n+1}(p,q)(Z))
=
g_n(p, q)(\del_{n+1} Z)
=
g_n(p,q)(X+Y)
=
(X,Y)
.
\]

We want to show that
\[
|Z| \geq |Z_X| + |Z_Y|
\]
because then we'll be done:
\[
\Zvol(X+Y) = |Z| \geq |Z_X| + |Z_Y| \geq \Zvol(X) + \Zvol(Y)
\]

To prove that
$|Z| \geq |Z_X| + |Z_Y|$,
we'll show that under the map $g_{n+1}(p,q)$,
any $t \in Z$ dies either in $Z_X$ or in $Z_Y$.

We'll call any $2$-simplex $t \in Z$ a `tet', short for `tetrahedron'.

Say that a tet $t \in Z$ has type $CCXY$ if $t= \pm[c_1c_2x_1y_1]$
for $c_1,c_2 \in C$, $x_1 \in A \setminus C$, $y_1 \in B \setminus C$.
Similarly for types $XXXX$, $CXXX$, $CXXY$, $XXXY$, etc.
The first two are pure $X$ cases, meaning that they live in $C_3(A)$;
the last two are hybrid cases.

Any $XX..$ tet dies in $Z_Y$;
any $YY..$ tet dies in $Z_X$.

The remaining cases to check are
the hybrid case $CCXY$
and the pure cases $CCCX$, $CCCY$.
The more interesting case is $CCXY$:
The key is that since $|C|=3$, $\{c_1,c_2\}$ cannot be disjoint from $\{p,q\}$.
As for $CCCX$, these must die in $Z_Y$
because $q \in \{c_1,c_2,c_3\} = C$.
Ditto for $CCCY$.

So 
\[
|Z| = |Z_X(p,q)| + |Z_Y(p,q)|,
\]
and $\Zvol$ adds.

Note how we're now emphasizing the possible dependence of $Z_X,Z_Y$ on $p,q$.
When $n=1$ the choice of $p,q$ can indeed make a difference:
We can get a different pair $(Z_X,Z_Y)$ if we switch
$p$ and $q$.
(Think about the connected sum of two cycle graphs.)

But when $n \geq 2$ we will now show that all tets are pure $X$ or pure $Y$,
so $Z_X$ consists of all the pure $X$ tets, and ditto for $Z_Y$.
This will make
$Z = Z_X + Z_Y$ with $Z_X \in \taut(X)$, $Z_Y \in \taut(Y)$.

Again our test case is $n=2$.

The key observation is that, now that we know that every tet dies on
one side or the other, we know that no tet can die on both sides,
because that would make $|Z| > \Zvol(X) + \Zvol(Y)$.

An $XXYY$ tet would die on both sides, so there can be none of these.

Any $pqXY$, $pXXY$, or $qXYY$ would die on both sides, and $p,q$ are arbitrary,
so this rules out all $CCXY,CXXY,CXYY$.

The only remaining hybrids are $XXXY$ and $XYYY$.
Let's rule out $XXXY$, leaving $XYYY$ to symmetry.

Fix any $y_0 \in B \setminus C$, and suppose there is
some $XXXy_0$ tet in $Z$.
Let $U \subset Z$ consist of all such $XXXy_0$ tets.
We claim that $y_0 \notin \vertices(\del U)$.
For let $s$ be any 2-simplex of the form $XXy_0$,
the only kind of 2-simplex containing $y_0$ that could belong
to $\del U$.
Any $t \in Z$ of which $s$ or $-s$ is a face must be a hybrid tet,
and the only possibility is $XXXy_0$, so $t \in U$.
Because $\del Z = 0$ the signed multiplicity of $s$
vanishes in $\del Z$,
hence also in $\del U$.
So $y_0 \notin \vertices(\del U)$,
making $U$ a complete cone on $y_0$,
contradiction.
So there is no such $XXXy_0$ tet, ruling out $XXXY$,
and with it $XYYY$.

So there are no hybrid tets, meaning that $Z$ splits, as claimed.

That takes care of $n=2$.
Let's quickly look at $n=3$.
As $CCXY$ was the crux for $n=2$,
$CCCXY$ is the crux for $n=3$.
$pqCXY$ dies on both sides, and $p,q$ are arbitrary,
so $CCCXY$ dies.

Finally, let's consider what goes wrong when $n=1$.
Here the only hybrid type is $CXY$,
and neither $pXY$ nor $qXY$ dies on both sides.
so we can't rule this type out.
Failing that, the $XXy_0$ tets needn't form a complete cone,
because $[x_0x_1y_0)$ can continue across $[x_1 y_0]$ to $-[p x_1 y_0]$,
so we can't rule out $XXY$ either.
\proofend

The foregoing proof works equally well over $\QQ$, providing we
permit ourselves to work with fractional multisets.
Alternatively, we can clear denominators as in the introduction
above.
Either way, we have:

\begin{corollary} \label{cor1}
$\Qvol$ adds under almost disjoint union,
and for $n \geq 2$
taut $\QQ$-fillings split.
\proofend
\end{corollary}

\section{Triangulations}

We turn now to filling simplicial triangulations of $S^2$.
Let's begin by establishing terminology.

An \emph{$n$-simplex} $s$ is simply a set of size $n+1$.
Its \emph{faces} are its subsets,
which are $k$-simplices with $-1 \leq k \leq n$,
$-1$ being the dimension of the empty simplex.
A \emph{simplicial complex} $\sigma$ is a finite collection of simplices
closed under taking faces:
If $t \subset s \in \sigma$ then $t \in \sigma$.

For any $k$-simplex $s \in \sigma$ define the \emph{link}
\[
\link(s,\sigma) = \{t: t \cap s = \emptyset, s \cup t \in \sigma \}
.
\]
This is a simplicial complex, but
(except for $\link(\emptyset,\sigma)=\sigma$)
it is not a subcomplex of $\sigma$,
because we are taking simplices in $\sigma$ and knocking down
their dimension by $k+1$.

We are interested in particularly nice simplicial complexes
called normal pseudomanifolds, or as we will prefer to say,
\emph{clean $n$-complexes}.

\begin{enumerate}
\item
To start with,
a clean complex must be \emph{pure},
meaning that every simplex of $\sigma$ belongs to some $n$-simplex in $\sigma$.
To determine $\sigma$ we can specify its $n$-simplices,
and then throw in all their subsets.
So we can confound a pure complex with the set of its $n$-simplices.

\item
A clean complex is a \emph{pseudomanifold}:
Every $n$-simplex abuts at most one other $n$-simplex
across any given $n-1$-simplex.
Another way to say this is that
the link $\link(s,\sigma)$ of any $n-1$-simplex $s$
consists of either a single point 
(in which case $s \in \del \sigma$ is a \emph{boundary simplex}),
or two points
($s \in \sigma \setminus \del \sigma$ is an \emph{interior simplex}.)

\item
A clean complex is \emph{normal}:
For any simplex $s$ of dimension $0,\ldots,n-2$,
$\link(s,\sigma)$ is connected.
(Some definitions extend this requirement
to the empty simplex of dimension $-1$;
this forces $\sigma$ to be connected.)
Being normal means that $\sigma$ is just what you get by gluing its 
$n$-simplices along shared faces,
without extra identifications. This rules out, for example,
an icosahedron with a pair of opposite vertices identified.
\end{enumerate}

If $\sigma$ is clean, 
$\link(s,\sigma)$ is clean.
The boundary $\del \sigma$ of $\sigma$ is clean,
and closed:
$\del \del \sigma = \{\emptyset\}$.

If $M$ is an $n$-manifold,
say that a simplicial complex $\sigma$
is a \emph{simplicial triangulation} of $M$
if the geometric carrier of $\sigma$ is homeomorphic to $S^2$.
In this case $\sigma$ is necessarily a clean complex.
Secretly we imagine that we've prescribed a specific homeomorphism,
at least up to the point of picking out one particular orientation
$X(\sigma) \in C_n$ if $M$ is oriented,
but we don't insist upon this because nothing will depend on which
orientation we pick.

This notion of triangulation is not as general as you might want
for some purposes.
The double cover of a triangle is not a simplicial triangulation
of $S^2$,
because that would requires two distinct
$2$-simplices to have the same three edges.
Nor can any $2$-simplex have two of its edges
glued to one another.
Thurston
\cite{thurston:shapes}
allows such triangulations,
but it is unclear how important these are to
his theory of shapes of surfaces.
We don't allow them.

We continue to think of a chain $Y \in C_n$ as a multiset of
oriented $n$-simplices.
If this is only nominally a multiset (all multiplicites are $1$)
we'll call $Y$ \emph{simplicial},
and view it as a pure simplicial complex.
(To stickle, in doing this
we are viewing oriented simplices as a subclass of simplices,
and confounding a set of $n$-simplices with a pure $n$-complex.)
We'll call $Y$ \emph{clean} if it is simplicial
and the associated simplicial complex is clean.

\section{Filling a triangulation of the 2-sphere}

\begin{theorem} \label{th2}
Let $\sigma$ be a simplicial triangulation of $S^2$,
and $Y$ a taut filling of $\sigma$.
Then $Y$ is clean, and arises from
a simplicial triangulation of $B^3$.
\end{theorem}

\proofstart

Let $v,e,f$  count the vertices, edges, and faces of $\sigma$.
We have
$e=3v-6$, $f=2v-4$.
(Check: $3f=2e$, $v-e+f=2$.)

Consider a counterexample pair $(\sigma,Y)$ for which
$\Zvol(\sigma)=|Y|$ is minimal.
Obviously $v>4$.
We can assume that $\sigma$ is prime
(not a connected sum along a triangle),
and in particular (this is all we will need) that there is no
vertex of degree $3$.

Call a tet $t \in Y$ \emph{eligible} if it shares two faces 
with $\sigma$. (It can't share more, since $\sigma$ has no vertex of
degree $3$.)
Since
\[
|Y| = \Zvol(\sigma) \leq f-\maxdeg(\sigma) 
\]
there must be at least $\maxdeg(\sigma)$ disjointly eligible tets,
but we'll only need two:
One with faces $s_1,s_2 \in \sigma$,
the other with disjoint faces $s_3,s_4 \in \sigma$.
The $s_3,s_4$ tet can't also have $s_1$ or $s_2$ as a face,
as then there would be a degree-$3$  vertex in $\sigma$.
So for any face $s$ of $\sigma$
there is an eligible tet without $s$ as a face.

Now as indicated in the introduction above, 
when we remove an eligible tet $t$,
we get a taut filling $Y-t$ of
$\sigma^\prime$,
where $\sigma^\prime$ is either (1) a triangulation of $S^2$,
or (2) the almost disjoint union of
two triangulations $\sigma_1,\sigma_2$ of $S^2$
that are joined along an edge of $t$.
We claim $Y-t$ is simplicial:
In the case (1)
$Y-t$ is actually clean,
by minimality of $|Y|$.
In case (2)
$Y-t$ is a taut filling of an almost disjoint union,
and as such splits $Y-t=Y_1+Y_2$ where $Y_i$ is
a clean taut filling of $\sigma_i$.
In this case $Y-t$ isn't clean, because the link of the common
edge shared by $\sigma_1$ and $\sigma_2$ is disconnected,
but $Y-t$ is still simplicial.

The crux here is to show that $Y$ is clean.
For starters, $Y$ is simplicial.
The only way $Y$ could fail to be simplicial is if $t$ has multiplicity
$2$.
But then $Y-t^\prime$ wouldn't be simplicial for eligible $t^\prime$
distinct from $t$,
contradiction.

Now, then:

\begin{enumerate}

\item
$Y$ is a pseudomanifold.
For suppose some $2$-simplex $s$ occurs more than twice as the boundary
of a $3$-simplex of $Y$.
If $s \notin \sigma$,
it must occur at least twice plus and twice minus;
after removing any eligible tet it
still occurs either twice plus or twice minus in $Y-t$ in case (1),
or in one or the other of $Y_1,Y_2$ in case (2),
contradicting minimality.
If $s \in \sigma$,
it occurs at least twice plus in $Y$:
Remove some eligible tet that does not have $s$ as a face
to get a contradiction.

\item
$Y$ has no edge $\{a,b\}$ whose link
is disconnected.
For the link is a 1-dimensional complex whose edges correspond to
tets of $Y$ that contain $\{a,b\}$.
The only way a component of the link can have fewer than three
vertices is if consists of a single edge $\{c,d\}$, corresponding
to a tet $t=\{a,b,c,d\} \in Y$.
In this case $\{a,b\}$ must be an edge of $\sigma$,
and $\{a,b,c\},\{a,b,d\}$ the faces of $\sigma$
that are adjacent along $\{a,b\}$.
These are the only faces of $\sigma$ that contain $\{a,b\}$,
so no other component of the link can reduce to a single edge.
If the link is disconnected, removing any eligible tet
other than $t$ will
leave it disconnected, contradicting minimality.

\item
$Y$ has no vertex whose link is disconnected:
For if $\link(\{a\},Y)$ is not connected,
the only way removing a single tet $t$
can render the link connected
is if one component of the link has a single $2$-simplex
$\{b,c,d\}$ and $t=\{a,b,c,d\}$.
In this case $\{a,b,c\}$, $\{a,b,d\}$, $\{a,c,d\}$ must
all be faces of $\sigma$,
making $a$ a vertex of $\sigma$ of degree $3$, contradiction.

\item Hence $Y$ is clean.
\proofend

\end{enumerate}

Thanks to this theorem, we may take the liberty of saying that
a taut filling of $\sigma$ `is' a simplicial triangulation
$\tau$ of $B^3$.

\section{Shelling}

A \emph{shelling} of a simplicial triangulation $\tau$ of $B^3$
is an ordering 
$(t_1,\ldots,t_{|\tau|})$
of its tets such that
gluing on one at a time maintains a triangulation of $B^3$,
i.e., 
so that for $k=1,\ldots,|\tau|$ the subcomplex $\tau_k$ generated
by $\{t_1,\ldots,t_k\}$ is a simplicial triangulation of $B_3$.
Each new tet will be glued on along one face ($v$ increases by $1$);
two faces ($v$ stays the same);
or three faces ($v$ decreases by $1$).
Call the shelling \emph{monotone}
if you never glue along three faces, so that $v$ never decreases.
Call $\tau$ \emph{freely monotone shellable} if any $t \in \tau$
can serve as the initial tet of a monotone shelling.

\begin{theorem} \label{th3}
If $\sigma$ is a simplicial triangulation of $S^2$,
any taut filling $\tau$ of $\sigma$
is a freely monotone shellable
simplicial triangulation of $B^3$.
\end{theorem}

\proofstart
This is a corollary of the proof of Theorem \ref{th2}
that insists on being called a separate theorem.
\proofend

\section{Clique complex}

A simplicial complex is a \emph{clique complex} if every clique in
its $1$-skeleton occurs as a simplex.
For a simplicial triangulation $\sigma$ of $S^2$,
a pithy formulation is,
`Every face is a triangle; every triangle is a face.'
Whitney \cite{whitney:hamiltonian} identified this as the criterion for
$\sigma$ to have a hamiltonian circuit.

If every triangle is a face,
$\sigma$ is \emph{prime},
in that is not the connected sum of smaller complexes.
Otherwise we can decompose $\sigma$ as the connected sum of
prime components $\sigma_i$.
(Of course, to reconstruct $\sigma$ we need more than this, we need to know
just how they are connected.)
From the point of taut fillings these components are independent:
Any taut filling $\tau$ of $\sigma$ is a union of taut fillings $\tau_i$.
That each of these is freely shellable implies the same of the union.

Now we want to show that, whether or not $\sigma$ itself is prime,
any taut filling of $\sigma$ is a clique complex.

\begin{theorem}
If $\sigma$ is a simplicial triangulation of $S^2$,
any taut filling $\tau$ of $\sigma$ is a clique complex.
\end{theorem}

\proofstart
The goal here is to show that neither of these
two taboo configurations occurs
in $\tau$:
\begin{enumerate}
\item
The three edges of triangle $\{A,B,C\}$ without the 2-simplex $ABC$.
\item
The six edges of the $K_4$ $\{A,B,C,D\}$ without the 3-simplex $ABCD$.
\end{enumerate}
The strategy will be to consider a minimal counterexample,
and find a modification that would produce a smaller counterexample.

If $\sigma$ has a vertex of degree 3,
it is a connected sum with one summand a single tetrahedron.
We can pull off the vertex from $\sigma$, while removing the associated
tet from $\tau$, to produce a smaller counterexample.

Since there is no vertex of degree 3.
there are at least four 3-simplices of $\tau$ that each meet
the boundary in two faces, and these pairs of faces are disjoint.
Removing any one of these tets performs an `edge flip' on the boundary
$\sigma$,
which either (a) remains a triangulation of the 2-sphere,
or (b) becomes an almost disjoint union where
two such triangulations are joined along an edge.
We will show that some choice for the flip does not remove the taboo
configuration.
In case (a) we have a smaller counterexample already.
In case (b) the key is that if a triangle or $K_4$ occurs
in the almost disjoint union, it occurs on one side or the other,
also giving a counterexample.

Configuration 1:
When we remove a tet from $\tau$ to flip an edge of $\sigma$,
the only edge that disappears from $\tau$ is the edge that gets flipped.
There are four choices of the edge to flip, and three edges we must
avoid flipping.

Configuration 2:
Having dealt with configuration 1,
we know that along with the six edges of $\{A,B,C,D\}$,
$\tau$ must contain all four of the faces $ABC,ABD,ACD,BCD$.
At most two of these can be boundary faces, because
$\sigma$ has no vertex of degree 3.
If we avoid flipping an edge of these at-most-two faces,
the $K_4$ will remain.
But we have at least four choices of which edge to flip,
involving four disjoint pairs of boundary faces,
so at least two choices will avoid removing any edge of the $K_4$.
\proofend

\bibliography{taut}
\bibliographystyle{amsplain}
\end{document}
